\section{Site resolution within a host, by chemistry and site}
\label{si:withinhost}

This section asks whether the ability of the networks to order the oxygen sites of one host depends on the
chemistry of the host or on the site (extension, not registered; both networks at \numWithinHostBreakdownBudget{}
training hosts). The rules (groups, thresholds, metrics and a host-cluster bootstrap paired across models and networks) were logged in \texttt{docs/deviations.md}
before any number was computed, and the analysis reads the existing test predictions without training anything
(\texttt{results/within\_host\_breakdown.json}; \texttt{docs/within\_host\_physics.md}). The measure is the
within-host residual error: for a host with at least two tested sites, the prediction and the DFT energy of each
site are taken relative to their host means, and the mean absolute difference of these relative values is the
error that remains once every host-level offset is removed. Its reference is the zero-skill value, the mean
absolute deviation of the DFT energies from their host mean, which is the within-host error of any predictor that
is constant within a host, such as the physics floor B0 or a network blind to the site. A network has
within-host skill in a group when its within-host error lies below the zero-skill value.

Each site's prediction is the seed ensemble averaged over the resamples in which its host was tested, so the
breakdown covers \numWithinHostBreakdownGroupsPooledAllNHosts{} hosts and \numWithinHostBreakdownGroupsPooledAllNSites{}
sites. This averaging differs from the registered estimator, which averages errors per resample: it lowers the
pooled site MAE of the originally registered network by
\numWithinHostBreakdownSanityISiteMaeVOneSiteAveragedMinusRecordS{}~eV for S and
\numWithinHostBreakdownSanityISiteMaeVOneSiteAveragedMinusRecordD{}~eV for D, so the values of this section are
comparable with each other but not with the registered learning curves.

Over all hosts, both networks resolve sites within a host, with within-host errors of
\numWithinHostBreakdownGroupsPooledAllVTwoSWithinHostMaePoint{}~eV (S) and
\numWithinHostBreakdownGroupsPooledAllVTwoDWithinHostMaePoint{}~eV (D) on the primary network (v2, post hoc)
against a zero-skill value of \numWithinHostBreakdownGroupsPooledAllVTwoSZeroSkillMaePoint{}~eV
(Table~\ref{tab:withinhost}). On v2 the descriptors leave the within-host error unchanged, S $-$ D =
\numWithinHostBreakdownGroupsPooledAllDifferencesSMinusDVTwoWithinHostMaePoint{}~eV
\ci{\numWithinHostBreakdownGroupsPooledAllDifferencesSMinusDVTwoWithinHostMaeCiNinetyFiveItemOne{}}{\numWithinHostBreakdownGroupsPooledAllDifferencesSMinusDVTwoWithinHostMaeCiNinetyFiveItemTwo{}};
on the originally registered network (v1) D has the larger within-host error, S $-$ D =
\numWithinHostBreakdownGroupsPooledAllDifferencesSMinusDVOneWithinHostMaePoint{}~eV
\ci{\numWithinHostBreakdownGroupsPooledAllDifferencesSMinusDVOneWithinHostMaeCiNinetyFiveItemOne{}}{\numWithinHostBreakdownGroupsPooledAllDifferencesSMinusDVOneWithinHostMaeCiNinetyFiveItemTwo{}},
so on v1 the network given the descriptors orders the sites of a host less well than the structure-only network.
The chemistry enters mainly through the spread of the target itself. Hosts whose cations have an empty $d$ shell
(\dzero{}) have a smaller site-to-site spread than hosts with a filled $d$ shell (\dten{}), with zero-skill values of
\numWithinHostBreakdownGroupsDClassDZeroVTwoSZeroSkillMaePoint{} and
\numWithinHostBreakdownGroupsDClassDTenVTwoSZeroSkillMaePoint{}~eV, and both networks keep within-host skill in both
classes. The v1 penalty of the descriptors is concentrated in the \dten{} hosts
(S $-$ D = \numWithinHostBreakdownGroupsDClassDTenDifferencesSMinusDVOneWithinHostMaePoint{}~eV
\ci{\numWithinHostBreakdownGroupsDClassDTenDifferencesSMinusDVOneWithinHostMaeCiNinetyFiveItemOne{}}{\numWithinHostBreakdownGroupsDClassDTenDifferencesSMinusDVOneWithinHostMaeCiNinetyFiveItemTwo{}})
and in the hosts that contain a cation with a high perturbed-host-state rate (Nb, W, Ti, V, Ta or Mo;
\numWithinHostBreakdownGroupsPhsPronePhsProneDifferencesSMinusDVOneWithinHostMaePoint{}~eV
\ci{\numWithinHostBreakdownGroupsPhsPronePhsProneDifferencesSMinusDVOneWithinHostMaeCiNinetyFiveItemOne{}}{\numWithinHostBreakdownGroupsPhsPronePhsProneDifferencesSMinusDVOneWithinHostMaeCiNinetyFiveItemTwo{}}).
Across the thirds of host stability the within-host error barely changes, while the site error of S on v2 rises
from \numWithinHostBreakdownGroupsStabilityTercileLowerVTwoSMaePoint{}~eV in the most stable third to
\numWithinHostBreakdownGroupsStabilityTercileUpperVTwoSMaePoint{}~eV in the least stable third. This pattern is
consistent with an error in the host-level offset rather than in the ordering of sites, a reading that is not
tested here. Among the thirds of the band gap, the smallest-gap third has the largest within-host errors of both
networks, and the errors at the most and least negative Bader charges of the vacancy oxygen within a host are
close to each other (Figure~\ref{fig:si_withinhost_groups}).

These are descriptive comparisons. The breakdown reports an interval for every group, model, network and metric
without correction for multiplicity, so some intervals exclude zero by chance alone, and an isolated interval that
excludes zero is weak evidence; groups with fewer hosts than the small-group threshold carry no interval. The
cation groups overlap, because a host with several cations is counted under each, and they are confounded with
each other and with structure type. Every group summary describes the filtered universe, that is, deep or
localised neutral vacancies only. Physical readings of these patterns are hypotheses
(\texttt{docs/within\_host\_physics.md}) and are not tested here.

\begin{table}[htbp]
\caption{Within-host residual MAE (eV) at \numWithinHostBreakdownBudget{} training hosts by group (extension;
both networks; v2 post hoc). Hosts and sites: those with at least two tested sites, on which the within-host error
is defined. Zero skill: the within-host error of any predictor constant within a host. Point values; the
two-sided intervals of the host-cluster bootstrap, and the per-cation groups of Figure~\ref{fig:si_withinhost_groups}a, are
in \texttt{results/within\_host\_breakdown.json} (\texttt{payload.groups}). The unassigned $d$ class is below the
small-group threshold and is descriptive only.}
\label{tab:withinhost}
\centering\small
% Generated by paper/si/tools/make_si_tables.py from results/within_host_breakdown.json (payload.groups). Do not edit.
\begin{tabular}{@{}lrrrrrrr@{}}
\toprule
 & & & & \pairhead{Primary (v2)} & \pairheadlast{Registered (v1)} \\
\whrules
Group & Hosts & Sites & Zero skill & S & D & S & D \\
\midrule
All tested hosts & \numWithinHostBreakdownGroupsPooledAllNHostsByMetricWithinHostMae{} & \numWithinHostBreakdownGroupsPooledAllNSitesWithin{} & \numWithinHostBreakdownGroupsPooledAllVTwoSZeroSkillMaePoint{} & \numWithinHostBreakdownGroupsPooledAllVTwoSWithinHostMaePoint{} & \numWithinHostBreakdownGroupsPooledAllVTwoDWithinHostMaePoint{} & \numWithinHostBreakdownGroupsPooledAllVOneSWithinHostMaePoint{} & \numWithinHostBreakdownGroupsPooledAllVOneDWithinHostMaePoint{} \\
\addlinespace
\whgroup{(b) $d$-electron class of the host} \\
$d$ class \dzero{} & \numWithinHostBreakdownGroupsDClassDZeroNHostsByMetricWithinHostMae{} & \numWithinHostBreakdownGroupsDClassDZeroNSitesWithin{} & \numWithinHostBreakdownGroupsDClassDZeroVTwoSZeroSkillMaePoint{} & \numWithinHostBreakdownGroupsDClassDZeroVTwoSWithinHostMaePoint{} & \numWithinHostBreakdownGroupsDClassDZeroVTwoDWithinHostMaePoint{} & \numWithinHostBreakdownGroupsDClassDZeroVOneSWithinHostMaePoint{} & \numWithinHostBreakdownGroupsDClassDZeroVOneDWithinHostMaePoint{} \\
$d$ class \dten{} & \numWithinHostBreakdownGroupsDClassDTenNHostsByMetricWithinHostMae{} & \numWithinHostBreakdownGroupsDClassDTenNSitesWithin{} & \numWithinHostBreakdownGroupsDClassDTenVTwoSZeroSkillMaePoint{} & \numWithinHostBreakdownGroupsDClassDTenVTwoSWithinHostMaePoint{} & \numWithinHostBreakdownGroupsDClassDTenVTwoDWithinHostMaePoint{} & \numWithinHostBreakdownGroupsDClassDTenVOneSWithinHostMaePoint{} & \numWithinHostBreakdownGroupsDClassDTenVOneDWithinHostMaePoint{} \\
$d$ class, other & \numWithinHostBreakdownGroupsDClassOtherNHostsByMetricWithinHostMae{} & \numWithinHostBreakdownGroupsDClassOtherNSitesWithin{} & \numWithinHostBreakdownGroupsDClassOtherVTwoSZeroSkillMaePoint{} & \numWithinHostBreakdownGroupsDClassOtherVTwoSWithinHostMaePoint{} & \numWithinHostBreakdownGroupsDClassOtherVTwoDWithinHostMaePoint{} & \numWithinHostBreakdownGroupsDClassOtherVOneSWithinHostMaePoint{} & \numWithinHostBreakdownGroupsDClassOtherVOneDWithinHostMaePoint{} \\
$d$ class, unassigned & \numWithinHostBreakdownGroupsDClassUnassignedNHostsByMetricWithinHostMae{} & \numWithinHostBreakdownGroupsDClassUnassignedNSitesWithin{} & \numWithinHostBreakdownGroupsDClassUnassignedVTwoSZeroSkillMaePoint{} & \numWithinHostBreakdownGroupsDClassUnassignedVTwoSWithinHostMaePoint{} & \numWithinHostBreakdownGroupsDClassUnassignedVTwoDWithinHostMaePoint{} & \numWithinHostBreakdownGroupsDClassUnassignedVOneSWithinHostMaePoint{} & \numWithinHostBreakdownGroupsDClassUnassignedVOneDWithinHostMaePoint{} \\
\addlinespace
\whgroup{(c) cations with high perturbed-host-state rates} \\
\begin{tabular}[t]{@{}l@{}}perturbed-host-state-prone\\cation present\end{tabular} & \numWithinHostBreakdownGroupsPhsPronePhsProneNHostsByMetricWithinHostMae{} & \numWithinHostBreakdownGroupsPhsPronePhsProneNSitesWithin{} & \numWithinHostBreakdownGroupsPhsPronePhsProneVTwoSZeroSkillMaePoint{} & \numWithinHostBreakdownGroupsPhsPronePhsProneVTwoSWithinHostMaePoint{} & \numWithinHostBreakdownGroupsPhsPronePhsProneVTwoDWithinHostMaePoint{} & \numWithinHostBreakdownGroupsPhsPronePhsProneVOneSWithinHostMaePoint{} & \numWithinHostBreakdownGroupsPhsPronePhsProneVOneDWithinHostMaePoint{} \\
\begin{tabular}[t]{@{}l@{}}no perturbed-host-state-\\prone cation\end{tabular} & \numWithinHostBreakdownGroupsPhsProneNotPhsProneNHostsByMetricWithinHostMae{} & \numWithinHostBreakdownGroupsPhsProneNotPhsProneNSitesWithin{} & \numWithinHostBreakdownGroupsPhsProneNotPhsProneVTwoSZeroSkillMaePoint{} & \numWithinHostBreakdownGroupsPhsProneNotPhsProneVTwoSWithinHostMaePoint{} & \numWithinHostBreakdownGroupsPhsProneNotPhsProneVTwoDWithinHostMaePoint{} & \numWithinHostBreakdownGroupsPhsProneNotPhsProneVOneSWithinHostMaePoint{} & \numWithinHostBreakdownGroupsPhsProneNotPhsProneVOneDWithinHostMaePoint{} \\
\addlinespace
\whgroup{(d) host stability (formation energy per atom)} \\
most stable third & \numWithinHostBreakdownGroupsStabilityTercileLowerNHostsByMetricWithinHostMae{} & \numWithinHostBreakdownGroupsStabilityTercileLowerNSitesWithin{} & \numWithinHostBreakdownGroupsStabilityTercileLowerVTwoSZeroSkillMaePoint{} & \numWithinHostBreakdownGroupsStabilityTercileLowerVTwoSWithinHostMaePoint{} & \numWithinHostBreakdownGroupsStabilityTercileLowerVTwoDWithinHostMaePoint{} & \numWithinHostBreakdownGroupsStabilityTercileLowerVOneSWithinHostMaePoint{} & \numWithinHostBreakdownGroupsStabilityTercileLowerVOneDWithinHostMaePoint{} \\
middle third & \numWithinHostBreakdownGroupsStabilityTercileMiddleNHostsByMetricWithinHostMae{} & \numWithinHostBreakdownGroupsStabilityTercileMiddleNSitesWithin{} & \numWithinHostBreakdownGroupsStabilityTercileMiddleVTwoSZeroSkillMaePoint{} & \numWithinHostBreakdownGroupsStabilityTercileMiddleVTwoSWithinHostMaePoint{} & \numWithinHostBreakdownGroupsStabilityTercileMiddleVTwoDWithinHostMaePoint{} & \numWithinHostBreakdownGroupsStabilityTercileMiddleVOneSWithinHostMaePoint{} & \numWithinHostBreakdownGroupsStabilityTercileMiddleVOneDWithinHostMaePoint{} \\
least stable third & \numWithinHostBreakdownGroupsStabilityTercileUpperNHostsByMetricWithinHostMae{} & \numWithinHostBreakdownGroupsStabilityTercileUpperNSitesWithin{} & \numWithinHostBreakdownGroupsStabilityTercileUpperVTwoSZeroSkillMaePoint{} & \numWithinHostBreakdownGroupsStabilityTercileUpperVTwoSWithinHostMaePoint{} & \numWithinHostBreakdownGroupsStabilityTercileUpperVTwoDWithinHostMaePoint{} & \numWithinHostBreakdownGroupsStabilityTercileUpperVOneSWithinHostMaePoint{} & \numWithinHostBreakdownGroupsStabilityTercileUpperVOneDWithinHostMaePoint{} \\
\addlinespace
\whgroup{(e) PBEsol(+$U$) band gap} \\
smallest-gap third & \numWithinHostBreakdownGroupsBandGapTercileLowerNHostsByMetricWithinHostMae{} & \numWithinHostBreakdownGroupsBandGapTercileLowerNSitesWithin{} & \numWithinHostBreakdownGroupsBandGapTercileLowerVTwoSZeroSkillMaePoint{} & \numWithinHostBreakdownGroupsBandGapTercileLowerVTwoSWithinHostMaePoint{} & \numWithinHostBreakdownGroupsBandGapTercileLowerVTwoDWithinHostMaePoint{} & \numWithinHostBreakdownGroupsBandGapTercileLowerVOneSWithinHostMaePoint{} & \numWithinHostBreakdownGroupsBandGapTercileLowerVOneDWithinHostMaePoint{} \\
middle third & \numWithinHostBreakdownGroupsBandGapTercileMiddleNHostsByMetricWithinHostMae{} & \numWithinHostBreakdownGroupsBandGapTercileMiddleNSitesWithin{} & \numWithinHostBreakdownGroupsBandGapTercileMiddleVTwoSZeroSkillMaePoint{} & \numWithinHostBreakdownGroupsBandGapTercileMiddleVTwoSWithinHostMaePoint{} & \numWithinHostBreakdownGroupsBandGapTercileMiddleVTwoDWithinHostMaePoint{} & \numWithinHostBreakdownGroupsBandGapTercileMiddleVOneSWithinHostMaePoint{} & \numWithinHostBreakdownGroupsBandGapTercileMiddleVOneDWithinHostMaePoint{} \\
largest-gap third & \numWithinHostBreakdownGroupsBandGapTercileUpperNHostsByMetricWithinHostMae{} & \numWithinHostBreakdownGroupsBandGapTercileUpperNSitesWithin{} & \numWithinHostBreakdownGroupsBandGapTercileUpperVTwoSZeroSkillMaePoint{} & \numWithinHostBreakdownGroupsBandGapTercileUpperVTwoSWithinHostMaePoint{} & \numWithinHostBreakdownGroupsBandGapTercileUpperVTwoDWithinHostMaePoint{} & \numWithinHostBreakdownGroupsBandGapTercileUpperVOneSWithinHostMaePoint{} & \numWithinHostBreakdownGroupsBandGapTercileUpperVOneDWithinHostMaePoint{} \\
\addlinespace
\whgroup{(f) Bader charge of the vacancy oxygen} \\
most negative O charge & \numWithinHostBreakdownGroupsBaderRankLowestNHostsByMetricWithinHostMae{} & \numWithinHostBreakdownGroupsBaderRankLowestNSitesWithin{} & \numWithinHostBreakdownGroupsBaderRankLowestVTwoSZeroSkillMaePoint{} & \numWithinHostBreakdownGroupsBaderRankLowestVTwoSWithinHostMaePoint{} & \numWithinHostBreakdownGroupsBaderRankLowestVTwoDWithinHostMaePoint{} & \numWithinHostBreakdownGroupsBaderRankLowestVOneSWithinHostMaePoint{} & \numWithinHostBreakdownGroupsBaderRankLowestVOneDWithinHostMaePoint{} \\
intermediate O charge & \numWithinHostBreakdownGroupsBaderRankMiddleNHostsByMetricWithinHostMae{} & \numWithinHostBreakdownGroupsBaderRankMiddleNSitesWithin{} & \numWithinHostBreakdownGroupsBaderRankMiddleVTwoSZeroSkillMaePoint{} & \numWithinHostBreakdownGroupsBaderRankMiddleVTwoSWithinHostMaePoint{} & \numWithinHostBreakdownGroupsBaderRankMiddleVTwoDWithinHostMaePoint{} & \numWithinHostBreakdownGroupsBaderRankMiddleVOneSWithinHostMaePoint{} & \numWithinHostBreakdownGroupsBaderRankMiddleVOneDWithinHostMaePoint{} \\
least negative O charge & \numWithinHostBreakdownGroupsBaderRankHighestNHostsByMetricWithinHostMae{} & \numWithinHostBreakdownGroupsBaderRankHighestNSitesWithin{} & \numWithinHostBreakdownGroupsBaderRankHighestVTwoSZeroSkillMaePoint{} & \numWithinHostBreakdownGroupsBaderRankHighestVTwoSWithinHostMaePoint{} & \numWithinHostBreakdownGroupsBaderRankHighestVTwoDWithinHostMaePoint{} & \numWithinHostBreakdownGroupsBaderRankHighestVOneSWithinHostMaePoint{} & \numWithinHostBreakdownGroupsBaderRankHighestVOneDWithinHostMaePoint{} \\
\bottomrule
\end{tabular}

\end{table}

\begin{figure}[htbp]
\centering
\sifig{fig_si_withinhost_groups}
\caption{Within-host residual MAE of S (blue) and D (green) at the largest budget (extension) by (\textbf{a}) cation,
(\textbf{b}) $d$-electron class, (\textbf{c}) perturbed-host-state-prone cations, (\textbf{d}) formation-energy
thirds, (\textbf{e}) band-gap thirds and (\textbf{f}) vacancy-oxygen Bader-charge rank. Grey ticks mark zero skill,
bars are two-sided host-cluster bootstrap intervals, filled circles and solid bars show the primary network (v2,
post hoc) and open circles and dashed bars the originally registered network (v1). Data:
\texttt{results/within\_host\_breakdown.json}.}
\label{fig:si_withinhost_groups}
\end{figure}
\FloatBarrier
